\documentclass[14pt]{extarticle}
\usepackage[a4paper, total={6in, 9in}]{geometry}
\usepackage{graphicx} % Required for inserting images
\usepackage[utf8]{inputenc}
\usepackage{CJKutf8}
\usepackage{amsmath}
\usepackage{amsthm}
\usepackage{amssymb}
\usepackage{amsfonts}
\usepackage{xcolor}
\usepackage{stmaryrd}
\usepackage{centernot}
\usepackage{tikz-cd}
\usepackage{hyperref}
\usepackage{enumitem}
\usepackage{xcolor}
\pagecolor{black}
\definecolor{PINK}{RGB}{222, 144, 219}
\color{PINK}
\setlength{\parindent}{0pt}


\newcommand{\B}[1]{\mathbb{#1}}
\newcommand{\C}[1]{\mathcal{#1}}
\newcommand{\Floor}[1]{\lfloor{#1}\rfloor}
\newcommand{\abs}[1]{\left|{#1}\right|}
\newcommand{\norm}[1]{\left|\left|{#1}\right|\right|}
\newcommand{\E}{\mathbb{E}}
\renewcommand{\S}{\mathbb{S}}
\newcommand{\F}{\mathbb{F}}
\newcommand{\R}{\mathbb{R}}
\newcommand{\V}{\mathbb{V}}
\newcommand{\N}{\mathbb{N}}
\newcommand{\W}{\mathbb{W}}
\newcommand{\U}{\mathbb{U}}
\renewcommand{\P}{\mathbb{P}}
\newcommand{\PD}[2]{\frac{\partial{#1}}{\partial{#2}}}
\newcommand{\PDT}[2]{\frac{\partial^2{#1}}{\partial{#2}^2}}
\newcommand{\LINE}[1]{\noindent\makebox[\linewidth]{\rule{8in}{#1pt}}}
\renewcommand{\v}{\mathrm{v}}
\renewcommand{\u}{\mathrm{u}}
\renewcommand{\b}{\mathrm{b}}
\newcommand{\w}{\mathrm{w}}
\renewcommand{\r}{\mathrm{r}}
\newcommand{\x}{\mathrm{x}}
\renewcommand{\a}{\mathrm{a}}
\newcommand{\e}{\mathrm{e}}
\newcommand{\SP}[1]{\left\langle #1 \right\rangle}
\newcommand{\iso}[1]{\overset{#1}{\cong}}
\newcommand{\rank}{\mathrm{rank}}
\newcommand{\im}{\mathrm{im}}
\newcommand{\OP}{\oplus}
\newcommand{\OT}{\otimes}
\newcommand{\conj}[1]{\overline{#1}}
\newcommand{\mat}[1]{\begin{bmatrix}#1\end{bmatrix}}
\newcommand{\china}[1]{\text{\begin{CJK*}{UTF8}{gbsn} #1 \end{CJK*}}}


\newcommand{\cat}{\mathrm{Cat}}
\newcommand{\ob}{\mathrm{ob}}
\newcommand{\mor}{\mathrm{mor}}
\renewcommand{\hom}{\mathrm{hom}}


\newtheorem{theorem}{Theorem}
\newtheorem{lemma}[theorem]{Lemma}

\theoremstyle{definition}
\newtheorem{definition}[theorem]{Definition}

\pagestyle{headings}
\allowdisplaybreaks

\numberwithin{equation}{section}

\begin{document}
\newgeometry{left=1cm,right=1cm,top=0.5cm,bottom=0.5cm}
\pagenumbering{gobble}
\[\text{Find }\lim_{n\to\infty}\E\left[\max_{\underset{2^n\le p\le2^{n+1}}{p\in\B{Z}}}\frac{1}{n+1}\sum_{i=0}^{n}B_{\Floor{p/2^i}}\right]\text{ where }B_k\sim U(\{0,1\})\text{. What if }B_k\sim U([0,1])?\]
\LINE{0.3}
\begin{proof}
Let $B_k$ be i.i.d. and $B_k\sim \text{Bernoulli}\left(\frac{1}{2}\right)$, the quantity
\[M_n=\max_{2^n\le{p}\le2^{n+1}}\frac{1}{n+1}\sum_{i=0}^{n}B_{\Floor{p/2^i}}\]
represents the maximum average of $n+1$ independent Bernoulli$\left(\frac{1}{2}\right)$ ariables along a root-to-leaf path in a binary tree.
Let $Z_n$ be the minimum number of zeros on any such path.
Thus, \[\sum_{i=0}^{n}B_{\Floor{p/2^i}}=n+1-Z_n\]
The sequence $Z_n$ satisfies the recursion
\[Z_{n+1}=(1-B_{\text{root}})+\min\left(Z_n^{(1)},Z_n^{(2)}\right)\]
where $Z_n^{(1)},Z_n^{(2)}$ are independent copies of $Z_n$.
Starting from $Z_1\sim\text{Bernoulli}\left(\frac{1}{2}\right)$, by monotone convergence of the distribution, $\E[Z_n]$ converges to a finite constant.
Thus, $\E[M_n]=n+1-\E[Z_n]=n+1-\C{O}(1)$ and
\[\lim_{n\to\infty}\frac{\E[M_n]}{n+1}=\lim_{n\to\infty}\frac{n+1-\C{O}(1)}{n+1}=1\]
Let $B_k$ be i.i.d. and $B_k\sim U([0,1])$.
Trivially, the increaments of $B_k$ have mean $\mu=\frac{1}{2}$ and finite variance.
The binary branching random walk with i.i.d. increments of mean $\mu$ satisfies the strong law of large numbers
\[\frac{1}{n}\max_{\text{paths}}\sum_{i=1}^{n+1}B_{\text{node}}\to\mu\qquad\text{a.s.}\]
Thus, \[\lim_{n\to\infty}\E\left[\max_{\underset{2^n\le p\le2^{n+1}}{p\in\B{Z}}}\frac{1}{n+1}\sum_{i=0}^{n}B_{\Floor{p/2^i}}\right]=\mu=\frac{1}{2}\]
\end{proof}
\end{document}